\section{BACKGROUND AND SIGNIFICANCE}\label{background-and-significance}

Medical concept normalization is the entity-linking task of mapping
in-text mentions of biomedical concepts like diagnoses, medications,
labs, and procedures to entries in a standardized ontology or
controlled vocabulary such as ICD-10-CM, LOINC, or RxNorm
\autocite{Miftahutdinov2021,Xu2022}. It sits at the foundation of
healthcare data interoperability. Errors propagate downstream: a
misassigned diagnosis code can shift a hospital's case-mix index,
mis-stratify a research cohort, or corrupt the training data of the
next generation of models. For most of the last decade, the state
of the art for automated coding came from supervised neural
classifiers trained on labeled clinical text
\autocite{Mullenbach2018,Huang2022}. The arrival of general-purpose
large language models has changed the deployment picture: foundation
models can in principle perform code prediction with no task-specific
fine-tuning. Two open questions motivated this work. Do the
fabrication findings reported on earlier model generations still
hold for the frontier models shipping today? And what does the cost
side of the deployment tradeoff look like for the configurations
that do work?

The literature has documented the fabrication problem but on
generations that are now dated. Soroush et al.\ (2024) reported
that GPT-4 achieved only 46\% exact match on ICD-9 and 34\% on
ICD-10, with a substantial fraction of errors being non-existent
codes rather than real-but-wrong assignments
\autocite{Soroush2024}. Lee and Lindsey documented that LLMs
systematically struggle to distinguish real codes from
plausibly-formatted fakes \autocite{Lee2024}. Both findings are an
instance of the broader medical-hallucination phenomenon
\autocite{Kim2025,Ji2023}, and both predate Anthropic's Opus 4.7,
OpenAI's GPT-5.5, and Google's Gemini 3 Flash Preview. Whether the
fabrication rates reported on GPT-4 generalize to the current
generation has not been re-measured.

Cost is largely absent from the published evaluations. The field's
response to hallucination has been architectural: retrieval-grounded
pipelines \autocite{DasBaksi2024,Sarvari2025}, agentic walkers
\autocite{Motzfeldt2025}, neuro-symbolic verifiers
\autocite{HybridCode2025}, and broader benchmark expansions
\autocite{Almeida2025,Singh2025,Li2025,Gershon2025}. Across this
body of work, hallucination is rarely treated as an explicit
outcome, and per-prediction cost is rarely reported at all. As
LLM-augmented coding moves toward production deployment, the
operationally relevant question is which (model, prompting mode)
combination achieves a useful accuracy at a tolerable cost. The
literature does not yet answer that question.

We present \textbf{Phantom Codes}, a reproducible benchmark with
three contributions. The first is empirical coverage of the most
recent frontier LLMs: 24 (model, prompting-mode) configurations
spanning Claude Opus 4.7, Sonnet 4.6, and Haiku 4.5; GPT-5.5 and
GPT-4o-mini; and Gemini 2.5 Pro, 2.5 Flash, and 3 Flash Preview.
Several of these models were released within weeks of the evaluation
date. The second is per-prediction cost as an outcome alongside
accuracy: every prediction in the matrix records token usage and
dollar cost at runtime, computed against the published per-token
rates of each provider, and we report cost-per-correct (USD per
exact-match outcome) at both the per-call and per-correct level.
The third is a set of methodological choices that make the first
two contributions interpretable: a six-way outcome taxonomy in
which fabrication of non-existent codes (hallucination) and
abstention (\texttt{no\_prediction}) are explicit, mutually
exclusive buckets; a within-model ablation across zero-shot,
constrained, and retrieval-augmented prompting modes that holds
the model and input fixed; and an abbreviation-stress mode (D4)
that replaces canonical display strings with clinical jargon
(T2DM, HTN, CKD-3) to expose where string-matching baselines
collapse. The headline matrix runs against Synthea-generated FHIR
Bundles \autocite{Walonoski2018}, a freely-redistributable synthetic
patient dataset that gives full reproducibility from a pinned
Synthea version and seed.
